\documentclass[11pt]{article}

\usepackage[a4paper,margin=2.4cm]{geometry}
\usepackage[T1]{fontenc}
\usepackage[utf8]{inputenc}
\usepackage{parskip}
\usepackage{xcolor}
\usepackage{booktabs}
\usepackage{enumitem}
\usepackage{fancyvrb}
\usepackage{mdframed}
\usepackage{titlesec}
\usepackage{array}
\usepackage{longtable}
\usepackage{hyperref}

\definecolor{ixaccent}{HTML}{0E7A64}
\definecolor{ixred}{HTML}{C0432B}
\definecolor{ixblue}{HTML}{2E6FA3}
\definecolor{ixdim}{HTML}{4A5851}
\definecolor{codebg}{HTML}{F2F5F0}
\definecolor{codeborder}{HTML}{C7D1C1}
\definecolor{notebg}{HTML}{F1E4C8}
\definecolor{noteborder}{HTML}{A66E14}

\hypersetup{
  colorlinks=true,
  linkcolor=ixaccent,
  urlcolor=ixaccent,
  pdftitle={Ixmucan\'e GeoPower -- Website Editing \& Deployment Guide},
  pdfauthor={Prepared with Claude Code}
}

\titleformat{\section}{\Large\bfseries\color{ixaccent}}{\thesection}{0.6em}{}
\titleformat{\subsection}{\large\bfseries\color{black}}{\thesubsection}{0.6em}{}
\titlespacing*{\section}{0pt}{1.4em}{0.6em}
\titlespacing*{\subsection}{0pt}{1.1em}{0.4em}

\newmdenv[
  backgroundcolor=codebg,
  linecolor=codeborder,
  linewidth=0.8pt,
  roundcorner=3pt,
  innerleftmargin=10pt,
  innerrightmargin=10pt,
  innertopmargin=6pt,
  innerbottommargin=6pt,
  skipabove=8pt,
  skipbelow=8pt,
  fontcolor=black
]{codebox}

\newmdenv[
  backgroundcolor=notebg,
  linecolor=noteborder,
  linewidth=0.8pt,
  roundcorner=3pt,
  innerleftmargin=10pt,
  innerrightmargin=10pt,
  innertopmargin=6pt,
  innerbottommargin=6pt,
  skipabove=8pt,
  skipbelow=8pt
]{notebox}

\newcommand{\code}[1]{\texttt{\small #1}}
\newcommand{\file}[1]{\texttt{\small\color{ixaccent}#1}}

\fvset{fontsize=\small, frame=none, xleftmargin=0pt}

\title{\vspace{-1.5cm}\textbf{\color{ixaccent}Ixmucan\'e GeoPower}\\[4pt]\Large Website Editing \& Deployment Guide}
\author{Prototype repository: \file{github.com/julimo80/Ixmucane-geopower}}
\date{\today}

\begin{document}
\maketitle
\vspace{-1.2cm}

\section*{Overview}
The site is an Astro + Tailwind static site. Almost all real content --- copy,
stats, service descriptions, case studies, team bios, photo/video captions ---
lives in plain JSON files under \file{src/data/}. You should not need to
touch the \file{.astro} component files (the layout/code) for day-to-day
content edits; only for structural changes.

\section{Running the Prototype}

The project needs Node.js $\geq$ 22 and ffmpeg (for video work). Both are
installed in a dedicated conda environment, \code{ixmucane-web}, kept separate
from your other project environments so nothing else on the machine is
affected.

\begin{codebox}
\begin{Verbatim}
# one-time per terminal session
conda activate ixmucane-web

cd ~/Ixmucane-geopower

npm run dev        # dev server at http://localhost:4321, hot-reloads on save
\end{Verbatim}
\end{codebox}

\begin{notebox}
\textbf{Note:} a dev server may already be running in the background from an
earlier session. If \code{npm run dev} complains the port is taken, either
reuse the one already running (just edit \& refresh the browser) or stop it
first with \code{pkill -f "astro dev"}.
\end{notebox}

Other commands, same environment:

\begin{codebox}
\begin{Verbatim}
npm run build       # production build to dist/ -- also catches JSON/syntax
                     # errors before they reach a live site
npm run preview     # serves that production build locally, for a final
                     # sanity check before deploying
\end{Verbatim}
\end{codebox}

To stop the dev server: \code{Ctrl+C} in its terminal, or \code{astro dev stop}
(the project's own convention, see \file{AGENTS.md}).

\section{Where to Edit Content}

Everything below is in \file{src/data/*.json}. Edit the text, save, and the
dev server updates the page automatically.

\begin{longtable}{@{}p{4.6cm}p{9.6cm}@{}}
\toprule
\textbf{File} & \textbf{Controls} \\
\midrule
\endhead
\file{site.json} & Site name, tagline, nav links, hero headline/lead, stat
  strip, contact heading/email, footer note. \\[4pt]
\file{pages.json} & The four link cards on the Home page (Services,
  Technology, Case Studies, Team, Contact overview text). \\[4pt]
\file{flagship.json} & The ``AI Favorability Mapping'' flagship service card
  and its map caption. \\[4pt]
\file{services.json} & The seven remaining numbered services (02--08). \\[4pt]
\file{technology.json} & Technology page: capability list, the 5-step
  process pipeline, and the \code{methodology.items[].figcap} /
  \code{description} fields -- the three figure captions (geological model,
  permeability cross-sections, ESP/GBOP pump sizing). \\[4pt]
\file{caseStudies.json} & The three case-study cards. Set
  \code{"draft": false} once a case study is client-confirmed; the yellow
  draft banner disappears automatically. \code{chartCaption} controls the
  chart caption on LOG 03. \\[4pt]
\file{team.json} & Team bios and the publications list. \\[4pt]
\file{fields.json} & The six Guatemalan fields plotted on the favorability
  map (name, status, note). \\[4pt]
\file{gallery.json} & \textbf{Home page photo gallery}: 8 entries, each with
  \code{key} (must match a filename in \file{src/assets/field/}),
  \code{caption}, and \code{note}. \\[4pt]
\file{fieldReel.json} & \textbf{Home page video collage}: 4 clips, each with
  \code{key}, \code{file} (filename in \file{public/videos/}), and
  \code{alt} (screen-reader text only, not shown visually). \\
\bottomrule
\end{longtable}

\section{Editing Photos and Videos}

\subsection{If you edit an existing image and keep the same filename}
\textbf{Yes --- the build will use your edited version automatically.} Astro
re-processes every image from \file{src/assets/} on each build/dev reload
and writes a fresh, content-hashed output file, so there is no manual
cache-busting step on the build side. The only thing that can look stale is
your \emph{browser's} own cache of the old picture -- do a hard refresh
(Ctrl+Shift+R) if you don't see the change.

\subsection{If you rename a file or switch to a different photo}
This depends on \emph{how} that image reaches the page:

\begin{itemize}[leftmargin=1.4em]
  \item \textbf{Home page gallery \& video collage} (\file{gallery.json} /
    \file{fieldReel.json}): fully key-driven, no code changes needed. Drop
    the new file into \file{src/assets/field/} (photos) or
    \file{public/videos/} (clips), then update \code{key} / \code{file} in
    the JSON to match the new filename.
  \item \textbf{Everything else} (Hero, Nav logo, Team photos, Technology
    figures, Case Study chart): these are imported by exact path at the top
    of the relevant \file{.astro} component (e.g.\
    \file{src/components/Team.astro} line 4:
    \code{import julianPhoto from '../assets/team/julian-xicara.webp'}). To
    swap one of these, either \textbf{overwrite the existing file in place}
    (simplest -- no code edit), or replace the file and edit that one
    \code{import} line to point at the new filename.
\end{itemize}

\begin{notebox}
\textbf{Rule of thumb:} same filename, no code touch, ever. Different
filename: check whether it is JSON-driven (edit JSON only) or
component-imported (edit one import line).
\end{notebox}

\subsection{Adding a brand-new photo to the gallery (not just replacing one)}
\begin{enumerate}[leftmargin=1.4em]
  \item Put the file in \file{src/assets/field/}, e.g.\ \file{new-shot.jpg}.
  \item Add an entry to \file{gallery.json}:
\begin{codebox}
\begin{Verbatim}
{ "key": "new-shot", "caption": "...", "note": "..." }
\end{Verbatim}
\end{codebox}
  \item Save -- it appears in the grid automatically (the component globs
    the folder by key, no import statement needed for gallery photos).
\end{enumerate}

\section{Verifying Your Changes}
\begin{enumerate}[leftmargin=1.4em]
  \item Keep \code{npm run dev} running and check the page in a browser at
    \url{http://localhost:4321}.
  \item Before considering something final, run \code{npm run build} once --
    it will fail loudly on a malformed JSON file or a broken image path,
    which the dev server sometimes lets slide silently.
\end{enumerate}

\section{Before Going Live}

\subsection{Content still marked as placeholder}
The footer currently says outright: \textit{``Prototype for review --- case
studies, team bios, publications and contact email are placeholders pending
confirmation, not live content.''} Concretely, that means:

\begin{itemize}[leftmargin=1.4em]
  \item \textbf{Photo/video captions} -- every entry in \file{gallery.json}
    and \file{fieldReel.json} has a \code{[Placeholder]} note asking you to
    confirm the field, well name, and year before publishing.
  \item \textbf{Case studies} -- LOG 01 and LOG 02 in \file{caseStudies.json}
    are marked \code{"draft": true} (real content, unconfirmed client). LOG
    03 (Tecuamburro--Infiernitos) is real and confirmed.
  \item \textbf{Publications} -- \file{team.json} has one placeholder entry
    inviting further publications/conference proceedings to be added.
  \item \textbf{Contact email} -- \file{site.json}'s
    \code{contact.email} (currently \code{hello@ixmucanegeopower.com}) needs
    to be a real, monitored mailbox before launch.
  \item Once everything above is confirmed, delete or rewrite
    \code{site.json}'s \code{footerNote} -- it exists specifically to flag
    this prototype status to reviewers.
\end{itemize}

\subsection{Deployment (turning this into a live, public website)}
The project is a fully static site (\code{npm run build} produces plain
HTML/CSS/JS/images in \file{dist/} -- no server or database needed), so
hosting is simple:

\begin{enumerate}[leftmargin=1.4em]
  \item \textbf{Pick a static host} and connect it to the GitHub repository
    (\file{github.com/julimo80/Ixmucane-geopower}). Any of these
    auto-detect Astro, run \code{npm run build}, and redeploy on every push
    to \code{main}:
    \begin{itemize}[leftmargin=1.2em]
      \item \href{https://vercel.com}{Vercel}
      \item \href{https://netlify.com}{Netlify}
      \item \href{https://pages.cloudflare.com}{Cloudflare Pages}
    \end{itemize}
    All three have a free tier sufficient for a marketing site like this
    one, and all three provision HTTPS/SSL automatically.
  \item \textbf{Domain name.} If \code{ixmucanegeopower.com} (or whichever
    domain you want) is not already registered, register it through any
    domain registrar. Point its DNS at the host chosen in step 1 --
    each host's dashboard gives you the exact records to add.
  \item \textbf{Contact form (optional upgrade).} The Contact page currently
    uses a plain \code{mailto:} link (no backend needed, works immediately).
    If you would rather have an in-page form, a service like
    \href{https://formspree.io}{Formspree} or the hosting platform's own
    form handling (e.g.\ Netlify Forms) can be wired in without standing up
    a custom backend.
  \item \textbf{Analytics (optional).} None is wired in yet. A
    privacy-friendly option such as Plausible, or Google Analytics if you
    prefer, can be added as a small script tag in
    \file{src/layouts/Base.astro}.
  \item \textbf{SEO basics (optional, nice-to-have).} The page already has a
    title/meta description per page. Not yet present: a
    \file{sitemap.xml}, a real Open Graph share image, and
    \file{robots.txt} -- worth adding once the domain is live, low effort
    via an Astro integration (\code{@astrojs/sitemap}).
\end{enumerate}

\section*{Quick Reference}
\begin{codebox}
\begin{Verbatim}
conda activate ixmucane-web        # once per terminal
cd ~/Ixmucane-geopower

npm run dev                         # edit + preview live
npm run build                       # verify a clean production build
npm run preview                     # preview that production build

# edit captions/copy:
src/data/gallery.json               # Home page photo captions
src/data/fieldReel.json             # Home page video alt text
src/data/technology.json            # Technology page figure captions
src/data/caseStudies.json           # Case study text + draft flags
src/data/team.json                  # Team bios + publications
src/data/site.json                  # Site-wide copy, contact email, footer
\end{Verbatim}
\end{codebox}

\end{document}
